\section{Retrieval-Side Clarification Changes Trajectory Quality in a Backbone-Dependent Way}
\label{sec:bar-charts}

Figure~\ref{fig:trajectory-quality-bars} compares three settings across five multimodal backbones: the unmodified \emph{Original} trajectory, the first-round \emph{Disturb} rewrite, and the \emph{Oracle} rewrite. We report two complementary metrics. \emph{RTE} measures trajectory-level retrieval efficiency, while \emph{LJ Score} captures the downstream answer-quality signal.

All bars are computed over the same benchmark universe. For \emph{Disturb} and \emph{Oracle}, samples that are not rewritten retain their \emph{Original} score, so each bar is a full-set mean under a fixed denominator rather than a subset-only average. This matters: the figure reflects the net intervention effect on the benchmark as a whole, not only the behavior of the rewritten subset.

\begin{table}[t]
    \centering
    \caption{Overall mean trajectory quality and answer quality across models.}
    \label{tab:trajectory-quality-overall}
    \begin{tabular}{lcccccc}
        \toprule
        Model & Orig. RTE & Dist. RTE & Orac. RTE & Orig. LJ & Dist. LJ & Orac. LJ \\
        \midrule
        Qwen2.5-vl-7B   & 0.3077 & 0.2893 & 0.3468 & 0.2847 & 0.2647 & 0.3287 \\
        Qwen3-vl-4B     & 0.4037 & 0.3206 & 0.4333 & 0.3887 & 0.3460 & 0.4287 \\
        Qwen3-vl-8B     & 0.3602 & 0.2883 & 0.3757 & 0.3707 & 0.3400 & 0.4073 \\
        Qwen3-vl-32B    & 0.3527 & 0.3007 & 0.3543 & 0.4067 & 0.3973 & 0.4207 \\
        InternVL3.5-8B  & 0.1429 & 0.0959 & 0.1799 & 0.1287 & 0.1087 & 0.1507 \\
        \bottomrule
    \end{tabular}
\end{table}

The first pattern is that retrieval-side clarification is not monotonic. On \texttt{Qwen2.5-vl-7B}, \emph{Oracle} increases both \emph{RTE} and \emph{LJ Score} relative to \emph{Original}, suggesting that explicit referential grounding can improve both process and outcome on a weaker backbone. The \texttt{Qwen3} family behaves differently. \emph{Disturb} consistently suppresses trajectory activity, while \emph{Oracle} often recovers part of that loss and, in several cases, improves answer quality as well. The direction and magnitude of the effect therefore depend on the backbone rather than on the intervention alone.

The second pattern is that process quality and answer quality are only partially coupled. In several settings, \emph{RTE} moves in one direction while \emph{LJ Score} moves in another, which means that trajectory efficiency should be interpreted as a process diagnostic rather than a proxy for correctness. A shorter or less active trajectory is not necessarily worse, and a more active trajectory is not automatically better. The bar charts therefore expose a genuine process--outcome decoupling: the same clarification can change the search behavior substantially without producing a commensurate change in final answer quality.

This decoupling becomes especially visible when comparing backbone families. The lower-capacity model, \texttt{Qwen2.5-vl-7B}, is relatively receptive to clarification in the sense that the intervention can increase retrieval activity and improve answer quality. By contrast, the \texttt{Qwen3} models tend to respond by reducing retrieval activity under disturbance, with oracle rewriting recovering only part of the lost behavior. The strongest backbone, \texttt{Qwen3-vl-32B}, is particularly informative: it is the most stable in final quality, yet it still shows a measurable reshaping of trajectory behavior under rewriting. In other words, capacity does not eliminate sensitivity to retrieval-side reformulation.

Task structure adds another layer to the story. Across all models, \emph{Subproblem Aggregation} consistently yields the highest \emph{RTE} and the strongest \emph{LJ Score}, reflecting a structurally retrieval-intensive regime. \emph{Comparison} and \emph{Multi-hop} are more fragile, especially under \emph{Disturb}, while \emph{Entity Recognition} is the task most likely to benefit from \emph{Oracle}. This task-level pattern is consistent with the intuition that explicit grounding is most useful when the identity of the referent is the dominant source of uncertainty.

Taken together, the figure supports a more general interpretation: referential ambiguity is not merely an error source to be removed, but a control variable that interacts with the retrieval policy already encoded in the backbone. The same clarification intervention can be activity-positive on one model family and activity-negative on another, while the resulting change in trajectory efficiency is only partially aligned with answer quality. This is why the intervention should be understood as a mechanism for changing controllability at the query--retriever interface, rather than as a universally monotonic accuracy booster.

% If you need a short take-home sentence for the main text:
% Retrieval-side clarification is backbone-dependent: it can increase retrieval activity and answer quality on some models, suppress activity on others, and its effect on trajectory efficiency is only partially aligned with final correctness.

